% ==================================================
% CHAPTER 2: REQUIREMENTS AND ANALYSIS
% ==================================================

\chapter{Requirements and Analysis}
\label{ch:requirements}

This chapter follows the structure of an IEEE~830 software requirements
specification: the product context is in Chapter~\ref{ch:overview}, and the
specific requirements, their models and their traceability are below.

\section{Functional Requirements}

Each requirement has a unique identifier used in the use cases, stories,
design and tests (Section~\ref{sec:trace}).

{\footnotesize
\begin{longtable}{>{\ttfamily}p{0.07\textwidth} >{\raggedright\arraybackslash}p{0.24\textwidth} >{\raggedright\arraybackslash}p{0.24\textwidth} >{\raggedright\arraybackslash}p{0.255\textwidth} p{0.045\textwidth}}
\toprule
\textrm{\textbf{ID}} & \textbf{Requirement} & \textbf{Inputs $\rightarrow$ outputs} & \textbf{Validation and errors} & \textbf{Pri.} \\
\midrule
\endhead
FR-01 & Users sign in with Google or email; sign-ins with the same verified email share one library. & Credentials at Cognito $\rightarrow$ access and refresh tokens; library id & RS256 signature, expiry, issuer, token type, client id; \texttt{401} for a bad token, \texttt{503} if Cognito is unreachable on a first sign-in & H \\
FR-02 & The extension saves the current page in one click, confirms it and offers undo. & URL, title, text, HTML, canonical URL $\rightarrow$ \texttt{saved}, \texttt{unchanged} or \texttt{skipped}, with the page & \texttt{http(s)} only (\texttt{400}); at least 50 words (\texttt{422}); size limits; 5{,}000 pages (\texttt{403}); \texttt{429} over the rate limit & H \\
FR-03 & Optional automatic capture saves only pages the user has read and that look like content, never private pages. & A page in a tab and reading signals $\rightarrow$ a saved page listed in ``Auto-saved today'', or nothing & Address rules, 30\,s in view or half scrolled, content signals; the server re-checks and answers \texttt{skipped} & H \\
FR-04 & Users can exclude sites from automatic capture and remove the day's automatic saves. & A site or page id $\rightarrow$ updated exclusion list; page deleted & Hosts normalised; a page already gone (\texttt{404}) counts as removed & M \\
FR-05 & YouTube videos are saved through their captions, description and channel. & Video URL with captions fetched in the tab $\rightarrow$ saved video & Valid video id and exact host; size limits; videos over two hours or without text are skipped & M \\
FR-06 & Addresses that differ only in tracking parameters, \texttt{www.} or fragments are one page; saving again updates it. & URL and canonical URL $\rightarrow$ \texttt{page\_id} & Canonical trusted only on the same site and not a parent of the page & H \\
FR-07 & Users search with a description, typed or spoken, and get up to ten pages with snippets. & Query (up to 500 characters) $\rightarrow$ ranked pages with title, site, time and snippet & Empty query \texttt{400}; \texttt{401}; \texttt{429}; an empty list when nothing is close & H \\
FR-08 & Queries that name a site, a time or a URL rank matching pages higher. & Query words $\rightarrow$ boosted ranking & Site words only from the user's own domains; dates in UTC; on failure, an ordinary search & M \\
FR-09 & Each result can be summarised on demand. & Query and page id $\rightarrow$ summary text & Page must be the caller's (\texttt{404}); 20\,s limit; \texttt{502}/\texttt{503} if Gemini fails or is not configured & M \\
FR-10 & Users browse their library newest first, delete pages with undo, and delete all their data. & Page id or ``delete all'' $\rightarrow$ library list; \texttt{204} & Owner only (\texttt{404}); 6\,s undo; on failure the page returns with an error message & H \\
FR-11 & Users report a problem from the website or the extension; it is emailed to the maintainer. & Message and context $\rightarrow$ \texttt{202}, then an email & Length and content cleaned (\texttt{422}); 5 per hour (\texttt{429}) & L \\
FR-12 & A public health endpoint reports the state of Solr, the model and Gemini. & -- $\rightarrow$ status of each dependency & Public; reports \texttt{degraded} instead of failing & L \\
\bottomrule
\caption{Functional requirements (priority: H high, M medium, L low)}
\label{tab:fr}
\end{longtable}
}

\section{Use Cases}

\subsection{Use Case Diagram}

\begin{figure}[H]
\centering
\begin{tikzpicture}[node distance=3mm]
  \newcommand{\actor}[3]{%
    \draw[thick] (#1,#2+0.45) circle (2.2mm);
    \draw[thick] (#1,#2+0.23) -- (#1,#2-0.4);
    \draw[thick] (#1-0.35,#2) -- (#1+0.35,#2);
    \draw[thick] (#1,#2-0.4) -- (#1-0.25,#2-0.8);
    \draw[thick] (#1,#2-0.4) -- (#1+0.25,#2-0.8);
    \node[font=\small] at (#1,#2-1.1) {#3};}
  \actor{0}{-3.1}{Reader}
  \actor{7.5}{-5.4}{Maintainer}
  \tikzset{usecase/.append style={minimum width=38mm, minimum height=8mm, font=\footnotesize}}
  \node[usecase] (signin)  at (-3.75, 0.0)  {Sign in (Google / email)};
  \node[usecase] (save)    at (-3.75,-1.25) {Save a page};
  \node[usecase] (auto)    at (-3.75,-2.5)  {Auto-capture pages};
  \node[usecase] (exclude) at (-3.75,-3.75) {Exclude a site};
  \node[usecase] (yt)      at (-3.75,-5.0)  {Save a YouTube video};
  \node[usecase] (undo)    at (-3.75,-6.25) {Undo / remove a save};
  \node[usecase] (search)  at (3.75, 0.0)  {Search by description};
  \node[usecase] (voice)   at (3.75,-1.25) {Search by voice};
  \node[usecase] (summ)    at (3.75,-2.5)  {Summarise a result};
  \node[usecase] (lib)     at (3.75,-3.75) {Browse the library};
  \node[usecase] (del)     at (3.75,-5.0)  {Delete a page / all data};
  \node[usecase] (report)  at (3.75,-6.25) {Report a problem};
  \begin{scope}[on background layer]
    \node[group, fill=gray!4, fit=(signin)(undo), inner sep=7pt] (extbox) {};
    \node[group, fill=gray!4, fit=(search)(report), inner sep=7pt] (webbox) {};
  \end{scope}
  \node[font=\small\bfseries, anchor=south] at (extbox.north) {Chrome extension};
  \node[font=\small\bfseries, anchor=south] at (webbox.north) {Web app};
  \foreach \u in {signin,save,auto,exclude,yt,undo}
    \draw (-0.35,-3.1) -- (\u.east);
  \foreach \u in {search,voice,summ,lib,del,report}
    \draw (0.35,-3.1) -- (\u.west);
  \draw (7.15,-5.4) -- (report.east);
  \draw[dashed,->] (voice) -- node[lbl, right] {$\ll$extend$\gg$} (search);
\end{tikzpicture}
\caption{Use case diagram: system overview}
\label{fig:usecases}
\end{figure}

The diagram shows the two actors and the twelve use cases, grouped by where
they happen. Voice search extends search: it only fills the search box.

\subsection{Use Case Templates}

\newcommand{\usecase}[8]{%
\begin{table}[H]
\centering
\footnotesize
\begin{tabularx}{\textwidth}{>{\bfseries}p{0.17\textwidth} L}
\toprule
\multicolumn{2}{l}{\normalsize\textbf{#1}} \\
\midrule
Actor & #2 \\
Preconditions & #3 \\
Trigger & #4 \\
Main flow & #5 \\
Alternative flows & #6 \\
Postconditions & #7 \\
Requirements & #8 \\
\bottomrule
\end{tabularx}
\end{table}}

\usecase{UC-1 Save a page}{Reader}{Signed in to the extension; the page is a web address.}%
{The reader presses \textbf{Remember this page}.}%
{1.~The extension reads the page's HTML, text and canonical URL. 2.~It sends them to \texttt{POST /memory}. 3.~The backend verifies the token, normalises the address, cleans and chunks the text, embeds it and stores it. 4.~The popup shows ``Saved just now'' with Undo.}%
{3a.~Same text as before: \texttt{unchanged}, nothing re-processed. 3b.~Under 50 words: refused with a message. 1a.~A YouTube video: captions are fetched in the tab instead.}%
{The page is in the reader's library and searchable within about a second.}{FR-02, FR-05, FR-06}

\usecase{UC-2 Auto-capture pages}{Reader (and the extension, acting for them)}{Automatic capture is on; the reader is signed in.}%
{A page finishes loading.}%
{1.~The extension checks the address rules and excluded sites. 2.~It waits until the page has been read. 3.~It checks the content signals. 4.~It sends the page; the server re-checks and saves it. 5.~The page appears in ``Auto-saved today''.}%
{1a, 2a, 3a, 4a.~Any check fails: nothing is saved. 4b.~Already saved: left unchanged. 5a.~The reader presses \textbf{Never this site}: the page is deleted and the site excluded.}%
{Only read, content-like, non-private pages are stored.}{FR-03, FR-04}

\usecase{UC-3 Search by description}{Reader}{Signed in to the web app; the library has pages.}%
{The reader types or speaks a query and presses Enter.}%
{1.~The web app sends \texttt{POST /search}. 2.~The backend runs keyword and vector search together, plus site and time searches if the query names them. 3.~It fuses the rankings, groups chunks into pages and drops weak matches. 4.~It returns up to ten pages with snippets.}%
{3a.~Nothing close: an empty list and a tip to describe the idea. 4a.~The reader presses \textbf{Summarise}: a summary of that page is shown (UC-3 extended by FR-09).}%
{Results are shown; nothing is changed.}{FR-07, FR-08, FR-09}

\usecase{UC-4 Delete a page}{Reader}{Signed in; the page is in the library.}%
{The reader presses the bin icon (or \texttt{Shift+Del}).}%
{1.~The page disappears and ``Page deleted $\cdot$ Undo'' is shown. 2.~After 6\,s the web app sends \texttt{DELETE /memory/\{id\}}. 3.~The backend deletes every chunk of the page.}%
{1a.~Undo within 6\,s: the page returns, nothing is sent. 2a.~The tab is closed or reloaded: the delete is sent at once and stays hidden after a reload. 3a.~The request fails: the page returns with an error message.}%
{The page is gone from the library and from search.}{FR-10}

\usecase{UC-5 Sign in}{Reader; Amazon Cognito}{The reader has a Google account or an email address.}%
{The reader presses \textbf{Get started} or \textbf{Sign in}.}%
{1.~The web app creates a PKCE verifier and challenge and redirects to Cognito. 2.~The reader signs in. 3.~Cognito redirects back with a code, exchanged for tokens. 4.~On the first API request the backend verifies the token and finds the library for this sign-in.}%
{4a.~A first sign-in with a verified email already in use: it joins that library. 4b.~Cognito unreachable on a first sign-in: \texttt{503}, try again. 4c.~A refresh is cancelled or throttled: the session is kept.}%
{The reader is signed in to their own library.}{FR-01}

\section{User Stories and Story Cards}

Requirements were planned as user stories, each with acceptance criteria
that became tests.

\begin{table}[H]
\centering
\small
\begin{tabularx}{\textwidth}{>{\ttfamily}p{0.07\textwidth} L >{\ttfamily}p{0.1\textwidth}}
\toprule
\textrm{\textbf{ID}} & \textbf{Story} & \textrm{\textbf{Reqs}} \\
\midrule
US-1 & As a reader, I want to save the page I'm on with one click, so that I can find it later. & FR-02 \\
US-2 & As a reader, I want pages I read to be saved automatically, so that I don't have to remember to save them. & FR-03 \\
US-3 & As a reader, I want my inbox, bank and private documents never to be captured, so that I can leave auto-capture on. & FR-03 \\
US-4 & As a reader, I want to search by describing a page, so that I find it without its title or address. & FR-07 \\
US-5 & As a reader, I want to say where or when I saw a page, so that ``the reddit post from last week'' finds it. & FR-08 \\
US-6 & As a reader, I want a short summary of a result, so that I know it's the right page before opening it. & FR-09 \\
US-7 & As a reader, I want to delete a page and undo a mistake, so that my library stays mine. & FR-10 \\
US-8 & As a reader, I want my Google and email sign-ins to share one library, so that I see the same pages either way. & FR-01 \\
US-9 & As a maintainer, I want problem reports by email with context, so that I can fix issues users hit. & FR-11 \\
\bottomrule
\end{tabularx}
\caption{User stories}
\label{tab:stories}
\end{table}

\begin{figure}[H]
\centering
\begin{tikzpicture}
  \tikzset{card/.style={draw, rounded corners=2pt, fill=yellow!6, text width=43mm,
    inner sep=4pt, font=\scriptsize, align=left, minimum height=44mm, anchor=north}}
  \newcommand{\storycard}[7]{%
    \node[card] at (#1,#2) {\textbf{#3}\hfill\textit{#4}\\[1mm]
      #5\\[1.5mm]\textbf{Acceptance}\\#6\\[1.5mm]\textbf{Priority / size:} #7};}
  \storycard{0}{0}{US-1}{Save}{As a reader, I want to save the page I'm on with one click.}%
    {$\bullet$ Saved within seconds, with Undo\\
     $\bullet$ Saving again updates one entry\\
     $\bullet$ Pages under 50 words get a clear refusal}{High / 3}
  \storycard{4.95}{0}{US-3}{Capture}{As a reader, I want my inbox, bank and private documents never to be captured.}%
    {$\bullet$ Email, banking, private apps, local and account pages refused\\
     $\bullet$ Password fields and \texttt{noindex} refused\\
     $\bullet$ Server refuses them even from an old extension}{High / 5}
  \storycard{9.9}{0}{US-4}{Search}{As a reader, I want to search by describing a page.}%
    {$\bullet$ A paraphrase finds the page\\
     $\bullet$ Exact names and codes still match\\
     $\bullet$ Weak matches are left out, not padded}{High / 8}
  \storycard{2.475}{-4.9}{US-5}{Search}{As a reader, I want to say where or when I saw a page.}%
    {$\bullet$ Site names come from my saved domains\\
     $\bullet$ ``yesterday'', ``last week'', ``in march'' work\\
     $\bullet$ Hints boost; they never hide a better page}{Medium / 8}
  \storycard{7.425}{-4.9}{US-7}{Library}{As a reader, I want to delete a page and undo a mistake.}%
    {$\bullet$ Page disappears at once; Undo for 6\,s\\
     $\bullet$ Closing the tab still deletes\\
     $\bullet$ A reload never shows it again}{High / 3}
\end{tikzpicture}
\caption{Story cards, with priority and relative size in story points}
\label{fig:storycards}
\end{figure}

\section{Workflow and Data Flow Models}

\subsection{Activity and Swimlane Diagrams}

Figure~\ref{fig:activity} follows one search from query to result, with the
parallel searches as a fork and join. Figure~\ref{fig:swimlane} follows an
automatic capture across the four parties involved; a decision that says
``no'' ends the flow there (flow-final node), so a page that fails any test
is never sent or stored.

\begin{figure}[H]
\centering
\begin{tikzpicture}[>={Stealth[length=1.8mm]}, action/.append style={text width=34mm}]
  \tikzset{bar/.style={fill=black, minimum height=1.2mm, inner sep=0pt}}
  \newcommand{\flowfinal}[1]{%
    \draw[fill=white] #1 circle (1.6mm);
    \draw #1 ++(45:1.6mm) -- ++(225:3.2mm);
    \draw #1 ++(135:1.6mm) -- ++(-45:3.2mm);}

  \node[startdot] (s) at (0,0) {};
  \node[action] (a1) at (0,-0.9)  {Type or speak a query};
  \node[action] (a2) at (0,-2.0)  {Send it with the access token};
  \node[adecide] (d1) at (0,-3.2) {token\\valid?};
  \node[action] (a3) at (0,-4.45) {Read site, time and URL hints in the query};
  \node[bar, minimum width=96mm] (fork) at (0,-5.45) {};
  \node[action] (b1) at (-3.3,-6.4) {Keyword search (BM25), top 50 chunks};
  \node[action] (b2) at (0,-6.4)    {Embed query; vector search (KNN), top 50};
  \node[action] (b3) at (3.3,-6.4)  {Search within the named site or time [if any]};
  \node[bar, minimum width=96mm] (join) at (0,-7.4) {};
  \node[action] (a4) at (0,-8.3)  {Fuse ranks (RRF); group chunks into pages};
  \node[action] (a5) at (0,-9.4)  {Drop weak matches (relevance floor)};
  \node[adecide] (d2) at (0,-10.6) {any page\\left?};
  \node[action] (a6) at (0,-11.85) {Pick snippets; show up to 10 pages};
  \node[adecide] (d3) at (0,-13.05) {summary\\wanted?};
  \node[action] (a7) at (3.6,-13.05) {Summarise that page with Gemini};
  \node[action] (a8) at (0,-14.3) {Open the page};
  \node[enddot] (e) at (0,-15.1) {};
  \node[action] (n1) at (-3.6,-10.6) {Suggest describing the idea instead};
  \node[action] (n2) at (-3.6,-3.2) {Show the sign-in screen};

  \draw[->] (s) -- (a1);
  \draw[->] (a1) -- (a2);
  \draw[->] (a2) -- (d1);
  \draw[->] (d1) -- node[tlabel, right] {yes} (a3);
  \draw[->] (d1) -- node[tlabel, above] {no} (n2);
  \draw[->] (a3) -- (fork);
  \foreach \b in {b1,b2,b3} { \draw[->] (fork.south -| \b) -- (\b); \draw[->] (\b) -- (join.north -| \b); }
  \draw[->] (join) -- (a4);
  \draw[->] (a4) -- (a5);
  \draw[->] (a5) -- (d2);
  \draw[->] (d2) -- node[tlabel, right] {yes} (a6);
  \draw[->] (d2) -- node[tlabel, above] {no} (n1);
  \draw[->] (a6) -- (d3);
  \draw[->] (d3) -- node[tlabel, above] {yes} (a7);
  \draw[->] (a7) |- (a8);
  \draw[->] (d3) -- node[tlabel, right] {no} (a8);
  \draw[->] (a8) -- (e);
  \draw[->] (n2) -- ++(0,-0.8) coordinate (f2); \flowfinal{($(f2)+(0,-0.16)$)}
  \draw[->] (n1) -- ++(0,-0.8) coordinate (f1); \flowfinal{($(f1)+(0,-0.16)$)}
\end{tikzpicture}

\caption{Activity diagram: a search}
\label{fig:activity}
\end{figure}

\begin{figure}[H]
\centering
\begin{tikzpicture}[>={Stealth[length=1.8mm]}, action/.append style={text width=27mm}]
  % lanes: x edges and headers
  \foreach \x in {0,3.1,7.0,10.8,13.6} \draw[gray!60] (\x,0.55) -- (\x,-15.5);
  \draw[gray!60] (0,0.55) -- (13.6,0.55);
  \draw[gray!60] (0,-0.05) -- (13.6,-0.05);
  \draw[gray!60] (0,-15.5) -- (13.6,-15.5);
  \foreach \x/\t in {1.55/Reader, 5.05/Chrome extension, 8.9/FastAPI backend, 12.2/Apache Solr}
    \node[font=\small\bfseries] at (\x,0.25) {\t};

  % flow-final node (circle with a cross)
  \newcommand{\flowfinal}[1]{%
    \draw[fill=white] #1 circle (1.6mm);
    \draw #1 ++(45:1.6mm) -- ++(225:3.2mm);
    \draw #1 ++(135:1.6mm) -- ++(-45:3.2mm);}

  % ---- nodes ----------------------------------------------------------------
  \node[startdot] (start) at (1.55,-0.55) {};
  \node[action] (a1) at (1.55,-1.35) {Open a page in Chrome};
  \node[action] (a2) at (4.75,-2.55) {Check: auto-capture on, signed in, site not excluded, address not private};
  \node[adecide] (d1) at (4.75,-4.05) {eligible?};
  \node[action] (a3) at (4.75,-5.35) {Watch the reading: 30\,s in view, or scrolled halfway};
  \node[action] (a0) at (1.55,-5.35) {Read the page};
  \node[adecide] (d2) at (4.75,-6.65) {read?};
  \node[action] (a4) at (4.75,-7.9) {Extract HTML, text and content signals};
  \node[adecide] (d3) at (4.75,-9.2) {content?};
  \node[adecide] (d5) at (4.75,-10.4) {where?};
  \node[action] (a10) at (4.75,-12.0) {Clean, chunk, embed on the device; store encrypted (IndexedDB)};
  \node[action] (a5) at (8.9,-10.4) {Verify token; re-check address and \texttt{noindex}};
  \node[adecide] (d4) at (8.9,-11.65) {allowed?};
  \node[action] (a6) at (8.9,-12.9) {Clean, chunk, embed (an existing page is kept)};
  \node[action] (a7) at (12.2,-12.9) {Store chunks and vectors};
  \node[action] (a8) at (4.75,-14.2) {Add to ``Auto-saved today''};
  \node[action] (a9) at (1.55,-14.2) {See it in the popup; may Remove it};
  \node[enddot] (end) at (1.55,-15.1) {};

  % ---- flows ----------------------------------------------------------------
  \draw[->] (start) -- (a1);
  \draw[->] (a1.south) |- (a2.west);
  \draw[->] (a2) -- (d1);
  \draw[->] (d1) -- node[tlabel, right] {yes} (a3);
  \draw[->, densely dashed] (a0) -- node[tlabel, above] {scrolls} (a3);
  \draw[->] (a3) -- (d2);
  \draw[->] (d2) -- node[tlabel, right] {yes} (a4);
  \draw[->] (a4) -- (d3);
  \draw[->] (d3) -- node[tlabel, right] {yes} (d5);
  \draw[->] (d5) -- node[tlabel, above] {server} (a5);
  \draw[->] (d5) -- node[tlabel, right] {device (default)} (a10);
  \draw[->] (a10) -- (a8);
  \draw[->] (a5) -- (d4);
  \draw[->] (d4) -- node[tlabel, right] {yes} (a6);
  \draw[->] (a6) -- (a7);
  \draw[->] (a7.south) |- node[tlabel, pos=0.75, above] {saved} (a8.east);
  \draw[->] (a8) -- (a9);
  \draw[->] (a9) -- (end);

  % "no" branches end the flow
  \foreach \d in {d1,d2,d3} {
    \draw[->] (\d.east) -- node[tlabel, above] {no} ++(1.25,0) coordinate (f\d);
    \flowfinal{($(f\d)+(0.16,0)$)}
  }
  \draw[->] (d4.east) -- node[tlabel, above] {no} ++(1.25,0) coordinate (fd4);
  \flowfinal{($(fd4)+(0.16,0)$)}
\end{tikzpicture}

\caption{Swimlane activity diagram: automatic capture of a page}
\label{fig:swimlane}
\end{figure}

\subsection{DFD Level 0}

The context diagram shows MemoryLane as one process with its external
entities and the data that crosses the boundary.

\begin{figure}[H]
\centering
\begin{tikzpicture}
  \node[proc, minimum size=28mm] (ml) at (0,0) {\textbf{0}\\[0.5mm]MemoryLane};
  \node[entity] (user) at (-5.4, 0)   {Reader};
  \node[entity] (cog)  at (5.0, 2.0)  {Amazon Cognito};
  \node[entity] (gem)  at (5.0, 0)    {Google Gemini};
  \node[entity] (ses)  at (5.0,-2.3)  {Maintainer\\\scriptsize (email via SES)};
  \node[entity] (yt)   at (0,  -3.3)  {YouTube};

  \draw[flow] ([yshift=3mm]user.east) -- node[lbl, above] {pages, queries, deletes} ([yshift=3mm]ml.west);
  \draw[flow] ([yshift=-3mm]ml.west) -- node[lbl, below] {results, snippets, summaries} ([yshift=-3mm]user.east);
  \draw[flow, <->] (ml) -- node[lbl, above, sloped, pos=0.45] {tokens, verified email} (cog.west);
  \draw[flow, <->] (ml) -- node[lbl, above, pos=0.45] {page text, summary} (gem.west);
  \draw[flow] (ml) -- node[lbl, below, sloped, pos=0.45] {problem reports} (ses.west);
  \draw[flow] (yt) -- node[lbl, left] {captions, via the user's browser} (ml);
\end{tikzpicture}

\caption{DFD level 0 (context diagram)}
\label{fig:dfd0}
\end{figure}

\subsection{DFD Level 1}

Level 1 splits the system into six processes and three data stores. It is
balanced with level 0: the reader's pages, queries and deletes go in, and
results, snippets and summaries come out; tokens flow with Cognito,
page text with Gemini and reports to SES. Every flow that touches D1 carries
the caller's verified library id.

\begin{figure}[H]
\centering
\begin{tikzpicture}[scale=0.8, transform shape, proc/.append style={minimum size=18mm, font=\scriptsize}]
  \newcommand{\dstore}[2]{\makebox[6.6mm][l]{#1}#2}
  \tikzset{wl/.style={lbl, fill=white, inner sep=1pt}}
  % external entities
  \node[entity] (user) at (-5.7, 0)   {Reader};
  \node[entity] (cog)  at (-5.7, 4.2) {Cognito};
  \node[entity] (gem)  at (8.9, 0)    {Gemini};
  \node[entity] (ses)  at (8.6,-4.2)  {SES};

  % processes
  \node[proc] (p1) at (-2.3, 4.2) {\textbf{1}\\Authenticate};
  \node[proc] (p2) at (-2.3, 2.1) {\textbf{2}\\Capture,\\save};
  \node[proc] (p3) at (-2.3, 0.0) {\textbf{3}\\Search};
  \node[proc] (p5) at (-2.3,-2.1) {\textbf{5}\\Manage\\library};
  \node[proc] (p6) at (-2.3,-4.2) {\textbf{6}\\Report a\\problem};
  \node[proc] (p4) at (6.1, 0)    {\textbf{4}\\Summarise};

  % data stores
  \node[store, minimum width=31mm] (d2) at (2.6, 4.2)  {\dstore{D2}{accounts.db}};
  \node[store, minimum width=31mm] (d1) at (2.6, 1.4)  {\dstore{D1}{Solr index}};
  \node[store, minimum width=31mm] (d4) at (2.6,-1.4)  {\dstore{D4}{device store}};
  \node[store, minimum width=31mm] (d3) at (2.6,-4.2)  {\dstore{D3}{reports.db}};

  % flows from the reader
  \foreach \p/\t in {p1/sign-in, p2/page, p3/query, p5/delete, p6/message}
    \draw[flow] (user.east) -- node[lbl, sloped, above, pos=0.5] {\t} (\p.west);

  \draw[flow, <->] (p1) -- node[lbl, above] {code, tokens} (cog);
  \draw[flow, <->] (p1) -- node[lbl, above] {library id} (d2);
  % capture writes to the device store (default) or to Solr
  \draw[flow] (p2.east) -- node[wl, pos=0.3] {chunks (server)} ([yshift=1.5mm]d1.west);
  \draw[flow] (p2.south east) -- node[wl, pos=0.78] {chunks (default)} ([yshift=1.5mm]d4.west);
  % search reads both stores
  \draw[flow] ([yshift=-1.5mm]d1.west) -- node[wl, pos=0.25] {ranked chunks} (p3.north east);
  \draw[flow] ([yshift=-1.5mm]d4.west) -- node[wl, pos=0.25] {ranked chunks} (p3.south east);
  % manage lists and deletes in both
  \draw[flow, <->] (p5.north east) -- node[wl, pos=0.75] {list, delete} ([xshift=6mm]d1.south west);
  \draw[flow, <->] (p5.east) -- node[wl, pos=0.45] {list, delete} ([xshift=6mm]d4.south west);
  % overflow and summaries
  \draw[flow] ([xshift=-5mm]d4.north east) -- node[lbl, left] {oldest past the limit} ([xshift=-5mm]d1.south east);
  \draw[flow] (d1.east) -- node[lbl, sloped, above] {page text} (p4);
  \draw[flow] (d4.east) -- node[lbl, sloped, below] {page text} (p4);
  \draw[flow, <->] (p4) -- node[lbl, above] {summary} (gem);
  \draw[flow] (p6) -- node[lbl, above] {report} (d3);
  \draw[flow] (d3) -- node[lbl, above] {email} (ses);
\end{tikzpicture}

\caption{DFD level 1}
\label{fig:dfd1}
\end{figure}

\subsection{DFD Level 2}

Level 2 decomposes process 3, Search. Its inputs and outputs match process 3
at level 1 (a query in, ranked results out, chunks read from D1).

\begin{figure}[H]
\centering
\begin{tikzpicture}[proc/.append style={minimum size=17mm, font=\scriptsize}]
  \newcommand{\dstore}[2]{\makebox[6.6mm][l]{#1}#2}
  \node[entity] (user) at (-6.4, 0) {Reader};

  \node[proc] (p31) at (-3.5, 0)   {\textbf{3.1}\\Read site,\\time, URL};
  \node[proc] (p32) at (-0.6, 1.45) {\textbf{3.2}\\Keyword\\search};
  \node[proc] (p33) at (-0.6,-1.45) {\textbf{3.3}\\Embed and\\vector search};
  \node[proc] (p34) at (2.2, 0)    {\textbf{3.4}\\Fuse ranks,\\group pages};
  \node[proc] (p35) at (4.9, 0)    {\textbf{3.5}\\Relevance\\floor};
  \node[proc] (p36) at (4.9,-2.6)  {\textbf{3.6}\\Snippets};

  \node[store, minimum width=31mm] (d1a) at (-0.6, 3.15) {\dstore{D1}{Solr index}};
  \node[store, minimum width=31mm] (d1b) at (-0.6,-3.15) {\dstore{D1}{Solr index}};

  \draw[flow] (user) -- node[lbl, above] {query} (p31);
  \draw[flow] (d1a.west) -| node[lbl, pos=0.75, left] {user's sites} (p31.north);
  \draw[flow] (p31) -- node[lbl, sloped, above] {text, filters} (p32);
  \draw[flow] (p31) -- node[lbl, sloped, below] {query} (p33);
  \draw[flow] (d1a) -- node[lbl, right] {BM25 hits} (p32);
  \draw[flow] (d1b) -- node[lbl, right] {nearest chunks} (p33);
  \draw[flow] (p32) -- node[lbl, sloped, above] {ranked list} (p34);
  \draw[flow] (p33) -- node[lbl, sloped, below] {ranked list} (p34);
  \draw[flow] (p34) -- node[lbl, above] {pages} (p35);
  \draw[flow] (d1a.east) -- node[lbl, sloped, above] {similarities} (p35.north);
  \draw[flow] (p35) -- node[lbl, right] {kept pages} (p36);
  \draw[flow] (p36.south) -- ++(0,-1.25) -| node[lbl, pos=0.25, below] {results with snippets} (user.south);
\end{tikzpicture}

\caption{DFD level 2 of process 3, Search. D1 is drawn twice to avoid crossing lines.}
\label{fig:dfd2}
\end{figure}

\section{External Interface Requirements}
\label{sec:interfaces}

\begin{table}[H]
\centering
\footnotesize
\begin{tabularx}{\textwidth}{>{\raggedright\arraybackslash}p{0.2\textwidth} >{\raggedright\arraybackslash}p{0.22\textwidth} L}
\toprule
\textbf{Interface} & \textbf{Protocol} & \textbf{Data exchanged} \\
\midrule
\multicolumn{3}{l}{\textit{User interfaces}} \\
Web app & HTTPS, any modern browser & Search box, results with snippets, library, account menu; keyboard shortcuts \\
Extension popup and options & Chrome extension UI & Save button and status, ``Auto-saved today'', auto-capture switch, excluded sites \\
\multicolumn{3}{l}{\textit{Software interfaces}} \\
MemoryLane REST API & HTTPS, JSON, \texttt{Authorization: Bearer} & Pages, queries, results, summaries, reports (Table~\ref{tab:api}) \\
Apache Solr & HTTP JSON, inside the server only & Chunk documents, BM25 and KNN queries, deletes \\
Amazon Cognito & OAuth 2.0 with PKCE; JWKS; userInfo & Authorization codes, tokens, signing keys, verified email \\
Google Gemini & HTTPS API & One page's text and the query; the summary \\
Amazon SES & SMTP on port 587 (STARTTLS) & Problem report emails \\
YouTube & HTTPS from the user's own tab & Caption track, title, channel, length, description \\
\multicolumn{3}{l}{\textit{Hardware and communication}} \\
Hardware & -- & No special hardware; a microphone for voice search \\
Communication & TLS 1.2+ via Caddy, HSTS & All traffic between browser and server is encrypted \\
\bottomrule
\end{tabularx}
\caption{External interfaces}
\label{tab:interfaces}
\end{table}

\section{Data Requirements}

The table states what must be stored and how it must be treated; the
implemented schema is in Section~\ref{sec:storage}.

\begin{table}[H]
\centering
\footnotesize
\begin{tabularx}{\textwidth}{>{\raggedright\arraybackslash}p{0.13\textwidth} L L L}
\toprule
\textbf{Entity} & \textbf{Content and format} & \textbf{Validation and integrity} & \textbf{Retention and privacy} \\
\midrule
Page & URL, title, site, save time, readable text, short description & \texttt{http(s)} only; $\geq$50 words; size limits; one entry per normalised URL per library & Until the user deletes it; visible only to its owner; never in logs \\
Passage (chunk) & Up to 150 words with a 384-number vector & Derived from the page; replaced when the page changes & Deleted with its page \\
Sign-in link & Cognito sign-in id, library id, verified email, last check & Only verified emails link sign-ins; re-checked daily & Removed by ``delete all my data'' \\
Problem report & Message and the context shown to the user & Length limits; cleaned text; stored before sending & Kept until sent; removed by ``delete all'' \\
Session & Access token, refresh token, email & Verified on every request & Access token in memory; refresh token in browser storage; cleared on sign-out \\
Extension settings & Auto-capture switch, excluded sites, today's automatic saves & Hosts normalised & In the browser only; automatic-save list keeps 24\,h \\
\bottomrule
\end{tabularx}
\caption{Data requirements}
\label{tab:datareq}
\end{table}

\section{Quality and Operational Requirements}

\begin{table}[H]
\centering
\footnotesize
\begin{tabularx}{\textwidth}{>{\ttfamily}p{0.06\textwidth} >{\raggedright\arraybackslash}p{0.14\textwidth} L L}
\toprule
\textrm{\textbf{ID}} & \textbf{Quality} & \textbf{Requirement (target)} & \textbf{How it is met / measured} \\
\midrule
QR-01 & Performance & Median search time under 200\,ms on the server's CPU & 76\,ms median, 93\,ms 95th percentile (Section~\ref{sec:evaluation}) \\
QR-02 & Relevance & MRR@10 on SciFact at least the prototype's 0.754 & 0.786; Hit@1 0.71 \\
QR-03 & Capacity & 5{,}000 pages per account on one 2-vCPU server & Page limit enforced; one EC2 instance \\
QR-04 & Reliability & No data lost when a save is retried or raced & Per-user write lock; content hash; idempotent deletes \\
QR-05 & Availability & Services restart after a crash or reboot; optional services fail softly & Restart policies, health checks; summaries and reports optional \\
QR-06 & Security & Every request except health authenticated; no access to another user's pages & PKCE, RS256 checks, user filter in the Solr client; isolation tests \\
QR-07 & Privacy & Private pages never auto-saved; no URLs or queries in logs & Two-stage capture rules; log filters; a test checks the logs \\
QR-08 & Usability & Every action reachable by keyboard; usable on phone widths & Shortcuts (\texttt{/}, arrows, \texttt{S}, \texttt{Shift+Del}); responsive layout \\
QR-09 & Accessibility & Controls labelled for assistive technology & Buttons with text or \texttt{aria-label}; visible focus \\
QR-10 & Maintainability & All tests pass on every change; decisions recorded & 566 tests in CI; 11 ADRs; documentation site \\
QR-11 & Portability & Runs wherever Docker runs; extension in any Chromium browser & Docker Compose; Manifest V3 \\
\bottomrule
\end{tabularx}
\caption{Quality and operational requirements}
\label{tab:nfr}
\end{table}

\section{Business Rules}

\begin{table}[H]
\centering
\footnotesize
\begin{tabularx}{\textwidth}{>{\ttfamily}p{0.06\textwidth} L}
\toprule
\textrm{\textbf{ID}} & \textbf{Rule} \\
\midrule
BR-01 & Every page belongs to exactly one library; pages are never shared between users. \\
BR-02 & A library belongs to the first verified sign-in for an email; later sign-ins with the same verified email use it; an unverified sign-in claims nothing. \\
BR-03 & A library holds at most 5{,}000 pages; saved pages can still be updated at the limit. \\
BR-04 & Automatic capture is off by default and never saves email, banking, social media, private apps, local addresses, account, login or payment pages. \\
BR-05 & An automatic save never rewrites a page already in the library; only a manual save updates it. \\
BR-06 & A delete becomes final when its 6-second undo window ends; ``delete all'' removes pages, sign-in links and reports. \\
BR-07 & Page text is sent to Gemini only when the user asks for that page's summary. \\
BR-08 & Requests are rate-limited per user (Table~\ref{tab:api}). \\
BR-09 & Logs never contain saved URLs or search text. \\
\bottomrule
\end{tabularx}
\caption{Business rules}
\label{tab:rules}
\end{table}

\section{Priorities, Acceptance and Traceability}
\label{sec:trace}

Priorities are given in Table~\ref{tab:fr}. A requirement is accepted when
its test cases pass (Chapter~\ref{ch:testing}). Table~\ref{tab:trace} links
each functional requirement to its use case, stories, design models and
tests.

\begin{table}[H]
\centering
\footnotesize
\begin{tabularx}{\textwidth}{>{\ttfamily}p{0.07\textwidth} >{\ttfamily}p{0.07\textwidth} >{\ttfamily}p{0.1\textwidth} L >{\ttfamily}p{0.16\textwidth}}
\toprule
\textrm{\textbf{FR}} & \textrm{\textbf{UC}} & \textrm{\textbf{US}} & \textbf{Design models} & \textrm{\textbf{Tests}} \\
\midrule
FR-01 & UC-5 & US-8 & Fig.~\ref{fig:seq-signin}, \ref{fig:state-session}, \ref{fig:class-client} & TC-11, TC-12, TC-18 \\
FR-02 & UC-1 & US-1 & Fig.~\ref{fig:seq-save}, \ref{fig:class-backend} & TC-01, TC-02 \\
FR-03 & UC-2 & US-2, US-3 & Fig.~\ref{fig:swimlane} & TC-04, TC-05 \\
FR-04 & UC-2 & US-3 & Fig.~\ref{fig:swimlane}, \ref{fig:extension} & TC-20 \\
FR-05 & UC-1 & -- & Fig.~\ref{fig:seq-save} & TC-13 \\
FR-06 & UC-1 & US-1 & Fig.~\ref{fig:class-backend} & TC-01, TC-03 \\
FR-07 & UC-3 & US-4 & Fig.~\ref{fig:activity}, \ref{fig:seq-search}, \ref{fig:collab-search}, \ref{fig:dfd2} & TC-06, TC-16 \\
FR-08 & UC-3 & US-5 & Fig.~\ref{fig:dfd2} (3.1), \ref{fig:class-backend} & TC-07 \\
FR-09 & UC-3 & US-6 & Fig.~\ref{fig:activity} & TC-14 \\
FR-10 & UC-4 & US-7 & Fig.~\ref{fig:collab-delete}, \ref{fig:state-page}, \ref{fig:class-client} & TC-08, TC-09, TC-10 \\
FR-11 & -- & US-9 & Fig.~\ref{fig:dfd1} (6) & TC-15 \\
FR-12 & -- & -- & Fig.~\ref{fig:components} & TC-19 \\
\bottomrule
\end{tabularx}
\caption{Traceability from requirements to use cases, stories, design and tests}
\label{tab:trace}
\end{table}

\paragraph{Changes and unresolved requirements.} Two features were built and
later withdrawn as requirements: a library map and ``Export my data''
(ADRs 0009 and 0010). Automatic capture was tightened after an outdated
extension saved an email inbox (ADR 0011). FR-05 has no user story of its own:
it was added from the team's own feature proposal of 30~August
(Section~\ref{sec:iterations}). No
requirement is unresolved; the open items are improvements, listed as future
work in Section~\ref{sec:future}.
